% TOP_PROBLEM: {"problemId": "problem.bunyakovsky-conjecture", "canonicalTitle": "Bunyakovsky conjecture", "releaseRank": 102, "primaryDomain": "number_theory_arithmetic_geometry", "primaryDomainLabel": "Number theory & arithmetic geometry", "displayStatus": "open_disputed_claim", "releaseStatus": "open_disputed_claim"}
% !TeX program = lualatex
% Definitions and sources retained from research_batch_101_103.tex; research is in GPT_6_Astra_Ultra.
\documentclass[11pt]{article}
\usepackage[margin=25mm]{geometry}
\usepackage{fontspec}
\setmainfont{Latin Modern Roman}
\usepackage{amsmath,amssymb,mathtools}
\usepackage{unicode-math}
\setmathfont{Latin Modern Math}
\usepackage{enumitem,needspace}
\usepackage{xurl,hyperref}
\hypersetup{colorlinks=true,urlcolor=blue}
\setcounter{secnumdepth}{1}
\setlength{\parindent}{0pt}
\setlength{\parskip}{5pt}
\setlength{\emergencystretch}{3em}
\setlist[itemize]{leftmargin=3em,itemsep=5pt,topsep=4pt,parsep=0pt}
\allowdisplaybreaks
\newcommand{\N}{\mathbb{N}}
\newcommand{\Z}{\mathbb{Z}}
\newcommand{\Q}{\mathbb{Q}}
\newcommand{\R}{\mathbb{R}}
\newcommand{\C}{\mathbb{C}}
\newcommand{\F}{\mathbb{F}}
\newcommand{\A}{\mathbb{A}}
\newcommand{\PP}{\mathbb P}
\newcommand{\E}{\mathbb{E}}
\newcommand{\eps}{\varepsilon}
\newcommand{\dd}{\,\mathrm d}
\newcommand{\sref}[2]{\hyperlink{source:#1:#2}{[S#2]}}
\newcommand{\eref}[2]{\hyperlink{extra:#1:#2}{[E#2]}}
% Turn the next switch off for a shorter reading copy. The quotations preserve
% every exactTarget field, including qualifications omitted from a short summary.
\newif\ifshowcatalogue
\showcataloguetrue
\newcommand{\cataloguescope}[1]{\ifshowcatalogue
 \paragraph{Exact scope in the supplied catalogue.}
 \begin{quote}\small #1\end{quote}\fi}
\begin{document}
\setcounter{section}{101}
% ================================================================
% problemId: problem.bunyakovsky-conjecture
% source releaseRank: 102; source releaseStatus: open_disputed_claim
% exactTarget SHA-256: 62e270d269e82f53d0bcc187ac56a1e341785d173040d79276b391c1c38a7b76
\clearpage
\section{\texorpdfstring{Bunyakovsky conjecture}{Bunyakovsky conjecture}}\label{102}
\subsection{Definitions and mathematical statement}
For a nonconstant integer polynomial $f$, a fixed prime divisor is a prime $p$ dividing $f(n)$ for every integer $n$. Suppose $f$ is irreducible over $\Q$, primitive, and has positive leading coefficient, with
\[
 \forall p\in\mathcal P\ \exists n\in\Z:\ p\nmid f(n).
\]
The assertion is that the set $\{n\in\Z_{>0}:f(n)\in\mathcal P\}$ is infinite. Positive leading coefficient ensures that signs cannot obstruct eventual positive prime values. No explicit count or density is required.
\par\smallskip\textit{Formulation and concept references:} \sref{102}{1}.
\cataloguescope{Prove or disprove that every irreducible polynomial f in Z[x] with positive leading coefficient and no fixed prime divisor takes prime values for infinitely many positive integer inputs.}
\subsection{Short English statement}
Must every irreducible positive-leading integer polynomial without a fixed prime divisor produce infinitely many primes?
\subsection{Sources}
\begin{itemize}\raggedright
\item[\textnormal{[S1]}]\hypertarget{source:102:1}{} Eric W. Weisstein with Ed Pegg Jr., Bouniakowsky Conjecture, MathWorld--A Wolfram Resource.
\url{https://mathworld.wolfram.com/BouniakowskyConjecture.html}.
{\small\textit{Catalogue formulation source.}}
\item[\textnormal{[E1]}]\hypertarget{extra:102:1}{}
R. J. Lemke Oliver, \emph{Almost-primes represented by quadratic polynomials},
Acta Arithmetica \textbf{151} (2012), no.~3, 241--261.
Theorem 1 (squarefree $P_2$ values for admissible quadratics), and Section 2.2,
Lemma 2 (the linear-sieve functions and bilinear remainder).
\url{https://lemkeoliver.github.io/papers/04-Quadratic.pdf}.
{\small\textit{Author's manuscript consulted 25 September 2026.}}
\item[\textnormal{[E2]}]\hypertarget{extra:102:2}{}
T. Browning, E. Sofos, and J. Ter\"av\"ainen,
\emph{Bateman--Horn, polynomial Chowla and the Hasse principle with
probability 1}, arXiv:2212.10373v2, 21 May 2026.
Theorems 1.2 and 1.5, including the coefficient-box and exceptional-set
quantifiers. \url{https://arxiv.org/html/2212.10373v2}.
{\small\textit{Consulted 25 September 2026; results averaged over
polynomials are distinguished from an assertion for every fixed polynomial.}}
\item[\textnormal{[E3]}]\hypertarget{extra:102:3}{}
L. Grimmelt and J. Merikoski,
\emph{On the greatest prime factor and uniform equidistribution of quadratic
polynomials}, arXiv:2505.00493 (2025), abstract and introduction,
the $n^{1.312}$ greatest-prime-factor bound for $n^2+1$.
\url{https://arxiv.org/abs/2505.00493}.
{\small\textit{Consulted 25 September 2026; a greatest-prime-factor result,
not a prime-values theorem.}}
% END RESEARCH ADDITION REFERENCES-102
\end{itemize}

\end{document}
